% Part: proof-theory
% Chapter: natural-deduction
% Section: grafting

\documentclass[../../../include/open-logic-section]{subfiles}

\begin{document}

\olfileid{pt}{nat}{gra}

\olsection{నిరూపణల అంటుకట్టడం \usetoken{P}{proof}}

\Log{N1c}, \Log{N1i}లలో \tetoken{నిరూపణలు}{!!^{proof}s} అనేవి ఉపపత్తులు~$!B$ నుంచి \tetoken{సూత్రాలు}{!!{formula}s}~$!A$కు గల \tetoken{నిరూపణలు}{!!{proof}s}. $!B$కే \tetoken{ఒక నిరూపణ}{!!a{proof}} ఉందనుకుందాం. అప్పుడు~$!B$ అనే ఉపపత్తిని ఉపయోగించకుండా~$!A$కు \tetoken{ఒక నిరూపణ}{!!a{proof}} పొందడానికి,~$!B$ యొక్క \tetoken{నిరూపణను}{!!{proof}}, $!B$ నుంచి~$!A$ యొక్క \tetoken{నిరూపణతో}{!!{proof}} ఎలా కలపాలి?

ఒక మార్గం: \Intro{\lif}ను వర్తింపజేసి~$!A$ యొక్క \tetoken{నిరూపణ}{!!{proof}} నుంచి~$!B$ ఉపపత్తిని తొలగించడం. అప్పుడు వచ్చిన~$!B \lif !A$ యొక్క \tetoken{నిరూపణను}{!!{proof}},~$!B$ యొక్క \tetoken{నిరూపణతో}{!!{proof}} కలిపి ఇది పొందవచ్చు:
\[
\AxiomC{$\Discharge{!B}{x}$}
\DeduceC{$!A$}
\DischargeRule{\Intro{\lif}}{x}
\UnaryInfC{$!B \lif !A$}
\AxiomC{}
\DeduceC{$!B$}
\RightLabel{\Elim{\lif}}
\BinaryInfC{$!A$}
\DisplayProof
\]
ఇది సూటిగా ఉన్నా కొంత అనవసరంగా అనిపిస్తుంది: వెంటనే తొలగించడానికి మాత్రమే~$\lif$ను ఎందుకు ప్రవేశపెట్టాలి? $!B \lif !A$ ద్వారా వెళ్లే ఈ ``పక్కదారి''ని నివారించగలగాలి. (ఇలాంటి పక్కదారులన్నింటినీ నివారించవచ్చనేదే సాధారణీకరణ ప్రమేయంలోని విషయం.)
\sourcecorrection{OLTEPTNATGRA-001}{ఈ నిరూపణ వృక్షంలోని నిగమనం కోసం మూలం $!B \lif !C$ అని ముద్రించింది. ఉపనిరూపణ $!A$ను ఇస్తుంది; తదుపరి వివరణకు అనుగుణంగా $!B \lif !A$గా సరిచేశాం.}

\Log{N1c}, \Log{N1i}లలో ఈ సందర్భంలో ఆ పక్కదారిని సులభంగానే నివారించవచ్చు:~$!B^x$ ఉపపత్తుల స్థానంలో~$!B$ యొక్క మొత్తం \tetoken{నిరూపణను}{!!{proof}} ఉంచవచ్చు. ఇతర నిరూపణల ఉపపత్తులపై \tetoken{నిరూపణలను}{!!{proof}s} ఇలా ``అంటుకట్టడం'' ఉపయోగకరం; అందుకే దానిని నిర్వచనంగా నమోదు చేద్దాం.

\begin{defn}
$\delta_1$ అనేది తెరిచి ఉన్న ఉపపత్తి~$!B^x$ నుంచి~$!A$కు గల \Log{N1c}-\tetoken{నిరూపణ}{!!{proof}} (లేదా \Log{N1i}-\tetoken{నిరూపణ}{!!{proof}}), $\delta_2$ అనేది~$!B$కు గల \Log{N1c}-\tetoken{నిరూపణ}{!!{proof}} (లేదా \Log{N1i}-\tetoken{నిరూపణ}{!!{proof}}) అనుకుందాం. అప్పుడు $\Subst{\delta_1}{\delta_2}{!B^x}$ అంటే~$\delta_1$లో విసర్జించబడని ప్రతి~$!B^x$ స్థానంలో~$\delta_2$ను ఉంచిన ఫలితం.
\end{defn}

$\delta_1$, $\delta_2$లలో ఒకే గుర్తు కలిగిన వేర్వేరు ఉపపత్తి \tetoken{సూత్రాలు}{!!{formula}s} లేకుండా,~$\delta_2$లోని ఏ తెరిచి ఉన్న ఉపపత్తిలోనూ~$\delta_1$ యొక్క ఐగెన్ చరరాశి లేకుండా ఉంటే, ఫలితమైన $\Subst{\delta_1}{\delta_2}{!B^x}$ సరైన \tetoken{నిరూపణ}{!!{proof}}; దాని తెరిచి ఉన్న ఉపపత్తులు~$\delta_1$లోని ప్రతిస్థాపిత ఉపపత్తులు మినహా మిగిలినవి, అలాగే~$\delta_2$లోని తెరిచి ఉన్న ఉపపత్తులు. \tetoken{నిరూపణలను}{!!{proof}s} అంటుకట్టేటప్పుడు ఈ షరతులు సాధారణంగా నెరవేరుతాయి. $\delta_1$, $\delta_2$లు ఒక పెద్ద \tetoken{నిరూపణ}{!!{proof}} యొక్క ఉపనిరూపణలైతే మొదటి షరతు నెరవేరుతుంది. రెండో షరతు నెరవేరకపోతే, అది నెరవేరేలా~$\delta_1$ యొక్క ఐగెన్ చరరాశులకు కొత్త పేర్లు పెట్టవచ్చు.
\sourcecorrection{OLTEPTNATGRA-004}{మూలం అంటుకట్టిన నిరూపణలో రెండు నిరూపణల తెరిచి ఉన్న ఉపపత్తులన్నీ మిగులుతాయని చెప్పింది. ప్రతిస్థాపిత ఉపపత్తులు ఇక తెరిచి ఉండవు; వాటిని మినహాయించాం.}

దీనికి అనురూపమైన నిర్వచనం \Log{N2c}, \Log{N2i}లకు వర్తిస్తుంది.

\begin{defn}
$\delta_1$ అనేది~$x:!B, \Gamma_1 \Sequent !A$కు గల \Log{N2c}-\tetoken{నిరూపణ}{!!{proof}} (లేదా \Log{N2i}-\tetoken{నిరూపణ}{!!{proof}}), $\delta_2$ అనేది~$\Gamma_2 \Sequent !B$కు గల \Log{N2c}-\tetoken{నిరూపణ}{!!{proof}} (లేదా \Log{N2i}-\tetoken{నిరూపణ}{!!{proof}}) అనుకుందాం. అప్పుడు $\Subst{\delta_1}{\delta_2}{x:!B}$ అంటే~$\delta_1$లోని ప్రతి ఆక్సియమ్ $x:!B \Sequent !B$ స్థానంలో~$\delta_2$ను ఉంచి, అటువంటి ఆక్సియమ్‌ల కింద ఉన్న ప్రతి సీక్వెంట్ సందర్భానికి~$\Gamma_2$ను చేర్చిన ఫలితం.
\end{defn}
\sourcecorrection{OLTEPTNATGRA-002}{ఈ నిర్వచనంలో $\delta_2$ను మూలం \Log{N1c}/\Log{N1i} నిరూపణగా పేర్కొంది. కానీ అది సీక్వెంట్‌కు గల నిరూపణ; ఈ సందర్భానికి అనుగుణంగా \Log{N2c}/\Log{N2i}గా సరిచేశాం.}

\begin{prop}
$\delta_1 \Proves x:!B, \Gamma_1 \Sequent
!A$, $\delta_2 \Proves \Gamma_2 \Sequent !B$ అయితే,~$\delta_1$ యొక్క ఏ ఐగెన్ చరరాశీ~$\Gamma_2$లో లేకపోయినపుడు $\Subst{\delta_1}{\delta_2}{x:!B} \Proves \Gamma_1, \Gamma_2 \Sequent !A$.
\end{prop}

\begin{proof}
  $\pheight{\delta_1}$పై ఆగమనంతో.
  \sourcecorrection{OLTEPTNATGRA-003}{ఆగమన ప్రమాణంలో మూలం నిర్వచించని $\delta$ను ఉపయోగించింది; ప్రతిపాదనలోని మొదటి నిరూపణ $\delta_1$గా సరిచేశాం.}
\end{proof}

\end{document}
